\section{Communion and conversion}
The first interpretation concerns belonging: how the numbered people and
the multinational multitude can describe the fullness of God's people,
and how Christ's members remain united across death. Bede's explanation
of Abraham's seed and Augustine's account of the departed at the altar
give distinct reasons for understanding the saints as members of one
Church. The prayers express the relationship between the saints
honored and the living who ask their intercession.

The second concerns the kind of life Christ gives: how his refreshment
agrees with costly obedience, hunger for justice, and persecution.
Chrysostom reasons from Christ's humility and from the bondage of sin;
Augustine explains the formation of desire, purity, and peace. The
Gradual's difficult promise and the Communion's repeated Beatitudes
receive more sustained attention in that account.

Both interpretations distinguish present participation from final
completion. Augustine's account of the Church includes faithful
members still needing conversion and departed members still awaiting
resurrection. His account of the Beatitudes admits present fulfillment
while preserving a future transformation beyond expression.
The differences among the sources remain particular: Bede permits
alternative identifications of the numbered company; Augustine and
Chrysostom give different referents to the inherited earth; inward
order and outward reconciliation are compatible emphases within
their explanations of peace.

The feast's final petitions allow both subjects to remain concrete.
The people who have honored the saints ask to rejoice in them always,
to be helped by their prayers, and to withstand human dangers.
They seek a holiness already displayed in Christ's members and
still needing to be formed in their own lives.
